\chapter{Preparation}

\section{Quantum Computing Preliminaries}

\subsection*{The Basics of Quantum Computing}

Classical hardware encodes information with bits, taking on binary values of either a 0 or 1. In the case of quantum computers, information can be stored using ``qubits'', allowing a qubit to be in a state of $\ket{0}$, $\ket{1}$ or in some linear combination $\alpha\ket{0} + \beta\ket{1}$ by leveraging quantum superposition. The values $\alpha$ and $\beta$ are complex numbers referred to as probability amplitudes, where $\lvert\alpha\rvert^2$ and $\lvert\beta\rvert^2$ give the probability of finding the qubit in the $\ket{0}$ and $\ket{1}$ states respectively, additionally satisfying the normalisation condition of $\lvert\alpha\rvert^2 + \lvert\beta\rvert^2 = 1$.



\section{}